\chapter{可信测量、Roofline 与性能计数器}

性能工程必须先建立测量纪律。第一册已经讲过输入规模、环境、编译选项、计时方式、热路径隔离、正确性校验和基础 \cmd{perf stat}。本章不重复这些基本规则，而是讨论第二册真正需要的进阶测量：如何把多核程序分解成计算、内存、分支、TLB、同步、调度和 I/O 假设；如何用 Roofline 判断上限；如何用性能计数器和对照实验排除错误解释。

\section{测量对象}

测量前先定义对象。是测单次操作延迟，还是测批量吞吐；是包含内存分配，还是只测热循环；是冷缓存，还是热缓存；是单线程，还是多线程；是平均值，还是尾延迟；是关心总吞吐，还是关心每个 worker 是否均衡。不同对象需要不同方法，不能混用结论。

常见错误包括：把初始化时间算进内核，把编译器优化掉的循环当作结果，把首次缺页当作算法成本，把输出打印放进计时区，把不同 CPU 频率状态下的结果直接比较。

第二册还要特别避免一种错误：只测总时间，不看扩展曲线。多线程程序的核心问题不是“某一次运行快不快”，而是从 1 个线程到多个线程时，吞吐如何变化。一个程序在 2 个线程时加速，在 8 个线程时停滞，在 16 个线程时变慢，这条曲线本身就是证据。它提示瓶颈可能从单核执行转移到了内存带宽、锁竞争、调度、NUMA 或 I/O。

\section{计时器和统计方法}

计时器应使用单调时钟，避免系统时间调整影响。短任务要重复执行足够多次，并防止编译器消除结果。报告结果时不要只给一次运行时间，应至少给出多次运行的最小值、中位数和波动范围。最小值常代表干扰较少的情况，中位数更接近日常情况，尾部异常则提示系统噪声或资源竞争。

预热同样重要。第一次运行可能包含页错误、动态链接、分支预测冷启动、缓存冷启动和 CPU 频率爬升。若目标是冷启动延迟，就应该保留这些成本；若目标是热循环吞吐，就应该把它们排除。测量对象不同，预热策略不同。

多核 benchmark 还要记录线程绑定和运行顺序。若每轮测试先跑单线程再跑多线程，CPU 温度、频率和缓存状态可能系统性变化。更稳妥的做法是随机化版本顺序，保留原始样本，并记录每轮线程数、绑定策略和输入位置。若差异小于系统波动，就不要写强结论。

\section{Roofline 模型}

Roofline 用算术强度连接计算峰值和内存带宽。算术强度是操作数与数据移动字节数的比值。低算术强度程序通常受内存带宽限制，高算术强度程序才可能接近计算峰值。这个模型不是精确预测器，而是帮助判断优化方向。

数组求和每读取一个整数只做一次加法，算术强度低；矩阵乘如果分块得当，每个元素可以被多次复用，算术强度高。若一个内核已经接近内存带宽上限，继续手写更复杂的算术指令可能没有意义；若内核远低于计算上限且数据复用充分，就要检查 SIMD、指令级并行和执行端口。

Roofline 的价值在于先给出“最多可能多快”的粗上限。假设一个内核每处理一个元素读取 8 字节、写入 8 字节，只做少量整数运算，那么它的算术强度很低。若平台可持续内存带宽是 50 GiB/s，那么不管循环写得多漂亮，只要每个元素必须从内存读写 16 字节，吞吐上限就被数据移动限制。反过来，若内核通过分块让同一批数据被重复使用，算术强度提高，才有资格讨论 FMA、SIMD 宽度和执行端口。

\begin{mentalmodel}
Roofline 不是为了画漂亮图，而是为了阻止错误优化：内存带宽瓶颈不要只改算术指令，计算端口瓶颈不要只改数据结构，同步瓶颈不要只改循环展开。
\end{mentalmodel}

\section{从字节数推导瓶颈}

Roofline 的关键是先估算字节数。以 \code{std::int32_t} 数组求和为例，每个元素读取 4 字节，做一次整数加法。如果数组远大于缓存，主要流量来自内存读取。若机器可持续内存带宽是 40 GiB/s，那么理想情况下每秒最多读取约 100 亿个整数。实际还要扣除循环开销、预取效率、TLB 和其他系统干扰。

这个估算能提前判断优化方向。如果实测已经接近带宽上限，就不该期待换一种加法写法获得数量级提升。若实测只有带宽上限的一小部分，就要检查访问是否连续、是否触发 TLB miss、是否有分支、是否被编译器向量化、是否被线程调度干扰。

字节数估算要区分“算法逻辑字节”和“实际流量”。逻辑上，一个原子计数器每次只更新 8 字节；硬件上，它可能让一整条 cache line 在核心之间来回迁移。逻辑上，一个对象只读取一个字段；硬件上，cache line 可能带入多个冷字段。逻辑上，\code{mmap} 文件像数组；系统上，第一次访问可能触发缺页和 page cache 读取。第二册的实验报告必须说明这些差异，否则字节数估算会过于乐观。

\section{性能计数器}

Linux 上常用 \cmd{perf stat} 查看 cycles、instructions、branches、branch-misses、cache-misses、page-faults、context-switches 等指标。更深入时可以用 \cmd{perf record} 和 \cmd{perf report} 定位热点。不同 CPU 的事件名称不同，实验报告要记录命令和平台。

计数器要组合解释。IPC 低可能因为 cache miss，也可能因为分支失败或同步等待。cache-misses 高不一定说明程序差，流式扫描大数组本来就会不断 miss，但可能充分利用带宽。context-switches 多可能说明线程过多、阻塞频繁或系统有干扰。

进阶测量更重要的是“假设驱动”。如果怀疑前端瓶颈，就比较指令缓存、分支、uops 供给和代码布局；如果怀疑内存带宽，就比较工作集大小、线程数、带宽工具和 LLC miss；如果怀疑 TLB，就改变页大小或访问页数；如果怀疑锁竞争，就看吞吐曲线、等待时间、futex 路径和临界区长度；如果怀疑 NUMA，就改变 first-touch 和绑定策略。不要一次性收集一大堆事件再事后找故事。

\section{Top-Down 思维}

现代性能分析常把 CPU 周期粗分为前端受限、错误预测受限、后端受限和有效退休。不同平台工具名称不同，但思路相同：先判断核心为什么没有持续退休有用指令，再进入更细分类。前端受限时检查取指、译码、代码布局和间接跳转；错误预测受限时检查分支输入分布；后端受限时继续区分执行端口、内存延迟、内存带宽和同步等待；有效退休比例高时，说明核心大部分时间确实在做有用工作，优化要回到算法或工作量本身。

这套思维可以避免把 IPC 低简单等同于“CPU 不行”。同样是 IPC 低，链表随机访问、锁等待、page fault、系统调用阻塞和分支乱跳的解决方向完全不同。报告里应该写“当前证据更支持哪一类受限”，而不是只列一个 IPC 数字。

\section{同步和调度的测量}

多线程程序经常慢在不退休用户态指令的地方。线程可能睡在 futex 上，可能被调度器切走，可能自旋等待 cache line，可能阻塞在 I/O，也可能因为 cgroup 配额用完而被节流。只看 cycles 和 instructions，容易忽略这些等待。

测同步要记录至少三类事实。第一，吞吐随线程数怎样变化；第二，等待在哪里发生，是用户态自旋、内核 futex、条件变量、I/O 还是队列为空；第三，竞争对象是什么，是一把锁、一个原子变量、一条 cache line、一个队列还是一个分片。若能使用 \cmd{perf sched}、\cmd{perf lock}、\cmd{perf c2c} 或 eBPF 工具，应记录工具版本和权限；若不能使用，也要用对照实验替代，例如拆分计数器、增加 padding、改变线程绑定和改变队列容量。

\section{多核扩展曲线}

一个高质量性能报告不只给“优化前后快了多少”，还要给扩展曲线。横轴是线程数或 worker 数，纵轴可以是吞吐、耗时、加速比、效率或尾延迟。理想线性加速很少出现，重要的是解释拐点。若 1 到 4 线程接近线性，8 线程开始变慢，可能是共享缓存、内存带宽、NUMA 或同步热点；若从 2 线程开始就没有提升，可能是任务粒度太小、串行部分太大或全局锁太重。

扩展曲线还要配合正确性校验。并行程序在高线程数下更快但偶尔错误，不是优化成功，而是暴露了竞争。每轮性能样本都应关联校验结果，校验失败的样本不能混入性能统计后继续分析速度。

\section{实验记录模板}

一份合格性能记录至少包含：机器型号、核心和 NUMA 拓扑、内存容量、编译器和编译选项、输入规模、数据类型、线程数、绑定策略、运行次数、计时范围、正确性校验、perf 命令和原始输出摘要。缺少这些信息，别人无法复现实验，你自己几周后也很难判断数字是否可靠。

性能结论要写成因果假设，而不是口号。例如“并行版本在线程数超过 8 后不再加速，perf 显示 LLC miss 和内存带宽接近平台上限，因此瓶颈更像内存带宽，而不是线程创建成本”。这样的结论可以继续被验证或推翻。

本册建议把报告分为五段。第一段写被测问题和正确性 oracle。第二段写机器、系统、编译和绑定环境。第三段写实验矩阵，包括输入规模、线程数、版本、运行轮次和顺序。第四段写原始结果摘要，不只给平均值，也给中位数、最小值和异常。第五段写归因假设和下一步验证。若证据不足，要明确写“当前只能排除某些解释，不能确认最终瓶颈”。

\begin{keyidea}
测量不是为了得到一个数字，而是为了排除错误解释，缩小瓶颈范围，并验证优化是否真的改变了瓶颈。
\end{keyidea}

\begin{labbox}
实验：对 Compute Systems Labs 的顺序归约和并行归约运行 \cmd{perf stat}。记录时间、cycles、instructions、IPC、cache-misses 和 context-switches。解释并行版本是否真的更快，如果没有，瓶颈可能在哪里。
\end{labbox}
